% fig_vl15_effpot.tex — Die Zentrifugalbarriere: wie aus V(r) und dem
% Drehimpulsterm das effektive eindimensionale Potential wird, und wie
% dessen Mulde mit ell wandert.
%
%   * geometry comes from the radial equation derived in the lecture,
%
%         ( -hbar^2/2m d^2/dr^2 + hbar^2 ell(ell+1)/(2 m r^2)
%                                + V(r) - E ) u(r) = 0,
%
%     i.e. from  V_eff(r) = V(r) + hbar^2 ell(ell+1)/(2 m r^2).
%     Nothing here is sketched by eye.
%
%   * the lecture leaves V(r) open, so the plate fixes the one choice it
%     announces for the next lecture, V(r) = -e^2/(4 pi eps0 r).  With
%     lengths in a0 = 4 pi eps0 hbar^2/(m e^2) and energies in
%     1 Ry = hbar^2/(2 m a0^2) = 13.6 eV the whole figure collapses to
%
%         V_eff(x) = ell(ell+1)/x^2 - 2/x,      x = r/a0,
%
%     which is what is plotted — one line of algebra, no free parameters,
%     and it puts a real order of magnitude on the axes.
%
%   * the two exact facts the picture is built around, both obtained by
%     setting dV_eff/dr = 0 for V = -alpha/r:
%
%         r_min = ell(ell+1) a0,     V_eff(r_min) = 1/2 V(r_min).
%
%     The second is the reason panel (b) carries the grey curve -1/x:
%     every minimum of the family sits on it, exactly.  It is drawn
%     from x = 2 on, i.e. starting at the first minimum it explains.  That is the kind
%     of statement a figure can make in one stroke and prose cannot.
%
%   * panel (a) uses ell = 1 and the level E = -1/4 Ry.  The turning
%     points then follow from  2/x^2 - 2/x = -1/4, i.e. x^2 - 8x + 8 = 0,
%     so x = 4 -+ 2 sqrt(2) = 1.1716 and 6.8284 — the arrow ends are
%     computed, not placed.
%
%   * colour encodes ell and nothing else, using the ordered ramp
%     `tpclevels` (ell = 0,1,2,3 -> spacecadet, cgblue, birrenorange,
%     birrenred), so the cgblue curve in (a) is the same state as the
%     cgblue curve in (b).  Entry 3 of the ramp, munsell, is skipped for
%     the reason the reference plate fig_vl11_kasten gives: beside cgblue
%     it is the same teal to the eye, so it would break the very code it
%     is meant to carry.  The two ingredients V(r) and the centrifugal
%     term are neutral grey and distinguished by dash pattern, because
%     they are not separate states; likewise the level E and the locus
%     1/2 V(r) are grey reference lines.
%
%   * the region the barrier closes off is hatched, not filled, so it
%     survives greyscale printing — as does the whole plate, since dash,
%     dot and solid already separate the grey lines.
%
%   * both y labels are rotated onto the left flank rather than sitting
%     at the top of the axis: the top belongs to the panel titles, and
%     `V_eff [Ry]` is too long to share it.
%
% Panel (a) how the barrier is built: attraction -2/x, repulsion 2/x^2,
%           and their sum, with the classically allowed radial window at
%           E = -1/4 Ry.  The point is that u(r) is kept away from the
%           origin by a *kinetic* term.
% Panel (b) the family ell = 0,1,2,3.  The well moves out like ell(ell+1)
%           and shallows like 1/[ell(ell+1)]: at ell = 3 it sits twelve
%           Bohr radii out and is only 1/12 Ry deep.

\begin{center}
\begin{tikzpicture}

% ─────────────────────────────────────────────────────────────
%  (a)  V(r) + centrifugal term = V_eff, for ell = 1
% ─────────────────────────────────────────────────────────────
\begin{axis}[
    name=build,
    tpcplot,
    width=6.4cm, height=6.3cm,
    grid=none,
    xmin=0, xmax=10, ymin=-1.0, ymax=0.95,
    enlarge x limits=false, enlarge y limits=false,
    xtick={0,2,4,6,8,10},
    ytick={-1,-0.5,0,0.5},
    yticklabels={$-1$,$-0.5$,$0$,$0.5$},
    xlabel={$r/a_0$},
    ylabel={$V\,[\mathrm{Ry}]$},
    every axis x label/.style={at={(ticklabel cs:0.5)},
                               anchor=near ticklabel, font=\small},
    every axis y label/.style={at={(axis description cs:-0.18,0.5)},
                               rotate=90, anchor=south, font=\small},
    domain=0.05:10, samples=500,
    restrict y to domain=-1.35:1.25,
    title={(a) $V_{\mathrm{eff}}$ \tpcifenglish{from two terms}%
                                {aus zwei Termen}, $\ell=1$},
]

  % The inner region closed off by the barrier at E = -1/4 Ry:
  % 0 < x < 4 - 2 sqrt(2) = 1.1716.  Two strips, because \closedcycle
  % closes to the line y = 0.
  \addplot[forbidden, draw=none, domain=0:1.1716, samples=2] {0.95}
    \closedcycle;
  \addplot[forbidden, draw=none, domain=0:1.1716, samples=2] {-1.0}
    \closedcycle;

  % Ingredient 1: the attraction, V(r) = -2/x  (grey, long dashes)
  \addplot[birrengray, line width=0.9pt, dash pattern=on 4pt off 2.5pt]
    {-2/x};

  % Ingredient 2: the centrifugal term, ell(ell+1)/x^2 = 2/x^2
  \addplot[classicallim, line width=1.0pt] {2/x^2};

  % Their sum: the effective potential for ell = 1
  \addplot[wavefn, cgblue] {2/x^2 - 2/x};

  % Minimum at x = ell(ell+1) = 2, depth -1/[ell(ell+1)] = -1/2
  \addplot[only marks, mark=*, mark size=1.8pt, cgblue]
    coordinates {(2,-0.5)};

  % A level in the well, E = -1/4 Ry: a neutral reference rule
  \addplot[levellinegrey, domain=0:10, samples=2] {-0.25};
  \node[lbl, anchor=north] at (axis cs:9.6,-0.28) {$E$};

  % The classically allowed radial window, x_- = 1.1716 to x_+ = 6.8284
  \draw[{Stealth[length=4pt]}-{Stealth[length=4pt]}, spacecadet,
        line width=0.7pt]
    (axis cs:1.1716,-0.25) -- (axis cs:6.8284,-0.25);
  \node[lbl, anchor=south] at (axis cs:4.0,-0.20)
    {\tpcifenglish{classically allowed}{klassisch erlaubt}};

  % The centrifugal term, named where it is steepest
  \node[lbl, text=birrengray, font=\scriptsize, anchor=north west,
        align=left] at (axis cs:2.95,0.68)
    {\tpcifenglish{centrifugal barrier}{Zentrifugalbarriere}\\[-1pt]
     $\dfrac{\hbar^2\ell(\ell+1)}{2mr^2}$};

  % The attraction: named and, once for the whole plate, defined.
  % A short leader ties the note to the dashed curve.
  \draw[birrengray, line width=0.4pt]
    (axis cs:6.0,-0.585) -- (axis cs:6.0,-0.338);
  \node[lbl, text=birrengray, font=\scriptsize, anchor=north west]
    at (axis cs:5.3,-0.60)
    {$V(r)=-\dfrac{e^2}{4\pi\varepsilon_0 r}$};

  % The sum, named just under its minimum
  \node[lbl, text=cgblue, anchor=north] at (axis cs:2.05,-0.53)
    {$V_{\mathrm{eff}}$};

  % Name the strip the barrier closes off
  \node[lbl, text=birrengray, font=\scriptsize, rotate=90,
        fill=white, fill opacity=0.85, text opacity=1, inner sep=1pt]
    at (axis cs:0.55,-0.50)
    {\tpcifenglish{classically forbidden}{klassisch verboten}};
\end{axis}

% ─────────────────────────────────────────────────────────────
%  (b)  The family V_eff(x) = ell(ell+1)/x^2 - 2/x
% ─────────────────────────────────────────────────────────────
\begin{axis}[
    tpcplot,
    at={(build.south east)}, anchor=south west, xshift=20mm,
    width=6.3cm, height=6.3cm,
    grid=none,
    xmin=0, xmax=18, ymin=-0.62, ymax=0.30,
    enlarge x limits=false, enlarge y limits=false,
    xtick={0,4,8,12,16},
    ytick={-0.6,-0.4,-0.2,0,0.2},
    yticklabels={$-0.6$,$-0.4$,$-0.2$,$0$,$0.2$},
    xlabel={$r/a_0$},
    ylabel={$V_{\mathrm{eff}}\,[\mathrm{Ry}]$},
    every axis x label/.style={at={(ticklabel cs:0.5)},
                               anchor=near ticklabel, font=\small},
    every axis y label/.style={at={(axis description cs:-0.20,0.5)},
                               rotate=90, anchor=south, font=\small},
    domain=0.05:18, samples=600,
    restrict y to domain=-0.80:0.45,
    title={(b) \tpcifenglish{The well moves out with $\ell$}%
                            {Die Mulde wandert mit $\ell$}},
]

  % The locus of the minima:  V_eff(r_min) = 1/2 V(r_min) = -1/x.
  \addplot[levellinegrey, domain=2:18, samples=200] {-1/x};

  % V_eff for ell = 0,1,2,3 — colour is the ordered ramp `tpclevels`
  \addplot[wavefn, spacecadet]   {-2/x};
  \addplot[wavefn, cgblue]       {2/x^2  - 2/x};
  \addplot[wavefn, birrenorange] {6/x^2  - 2/x};
  \addplot[wavefn, birrenred]    {12/x^2 - 2/x};

  % Their minima, at x = ell(ell+1) with depth -1/[ell(ell+1)]
  \addplot[only marks, mark=*, mark size=1.8pt, cgblue]
    coordinates {(2,-0.5)};
  \addplot[only marks, mark=*, mark size=1.8pt, birrenorange]
    coordinates {(6,-0.16667)};
  \addplot[only marks, mark=*, mark size=1.8pt, birrenred]
    coordinates {(12,-0.08333)};

  % Direct labels, no legend
  \node[lbl, text=cgblue, font=\scriptsize, anchor=north]
    at (axis cs:2.0,-0.53)  {$\ell=1$};
  \node[lbl, text=birrenorange, font=\scriptsize, anchor=south]
    at (axis cs:6,-0.135)   {$\ell=2$};
  \node[lbl, text=birrenred, font=\scriptsize, anchor=south]
    at (axis cs:12,-0.050)  {$\ell=3$};
  \node[lbl, text=spacecadet, font=\scriptsize, anchor=west]
    at (axis cs:4.3,-0.565) {$\ell=0$: $V(r)$};

  % The grey curve through the minima, with a short leader
  \draw[birrengray!70, line width=0.4pt]
    (axis cs:14.6,0.052) -- (axis cs:15.5,-0.0645);
  \node[lbl, text=birrengray, font=\scriptsize, anchor=south west]
    at (axis cs:13.9,0.055) {$\tfrac12 V(r)$};

  % The two exact statements the picture is built on
  \node[lbl, text=birrengray, font=\scriptsize, anchor=north west,
        align=left] at (axis cs:9.0,-0.30)
    {$r_{\min}=\ell(\ell+1)\,a_0$\\[1pt]
     $V_{\min}=\tfrac12 V(r_{\min})$};
\end{axis}

\end{tikzpicture}
\end{center}
